\documentclass{article}
\usepackage[utf8]{inputenc}
\usepackage{amsmath}
\usepackage{nd3} %% This is package for typesetting derivations.

\title{24.241 Pset Template}
\author{Your MIT ID}
\date{\today}

\begin{document}

\maketitle

Here is an example derivation:

\begin{ND}[][][][][0.7\linewidth]
    \ndl{1}{$P \rightarrow Q$}{A}
    \ndl{2}{$P$}{A}
    \ndl{1,2}{$Q$}{1, 2 MP}
\end{ND}

It is produced with the following code:

\begin{verbatim}
\begin{ND}[][][][][0.7\linewidth]
    \ndl{1}{$P \rightarrow Q$}{A}
    \ndl{2}{$P$}{A}
    \ndl{1,2}{$Q$}{1, 2 MP}
\end{ND}
\end{verbatim}

The lines of the derivation are enclosed by the \verb|ND| environment: \verb|\begin{ND}|\ldots\verb|\end{ND}|. (The \verb|[][][][][0.7\linewidth]| part just sets the sizing of the derivation; it's optional.)

To typeset a line, use \verb|\ndl{|lines of the assumptions\verb|}{|sentence\verb|}{|rule of inference\verb|}|.

Mathmatical symbols need to be enclosed by \verb|$|s, e.g., \verb|$P \rightarrow Q$| for $P \rightarrow Q$.

Here are some of the symbols we'll use:

\begin{itemize}
    \item $\neg$: \verb|\neg|
    \item $\land$: \verb|\land|
    \item $\lor$: \verb|\lor|
    \item $\rightarrow$: \verb|\rightarrow|
    \item $\forall$: \verb|\forall|
    \item $\exists$: \verb|\exists|
\end{itemize}

\end{document}
